% \iffalse meta-comment
%
% hit-key-engine.dtx 是各个类共用的键机制：\hitsetup 入口、键族、常量与字段的声明
%
% \fi
%
% \section{键机制（engine）}
%
% \changes{v3.2a}{2026/08/07}{\cs{hitsetup} 的底层机制抽出来给各个类共用}
% 各个类的元信息字段名单各不相同，但生成字段、拆关键词、分族路由、旧名兼容这套
% 机制是一样的，原先每个类各写一遍，改一处容易漏掉别处。这里只放机制，字段名单
% 与各族的键仍旧写在各个类里。
%
% 这个文件里全是定义，不执行任何依赖类状态的东西，所以排在类自己的模块之前也没
% 关系。唯一要各个类自己发动的是 \cs{\_\_hit\_check\_class\_options:}，它要读
% \cs{@classoptionslist}，得等旧名映射表填好之后再调用。
%
%    \begin{macrocode}
%<*shared-key-engine>
%    \end{macrocode}
%
% \begin{macro}{\hit@declare@class@key,\hit@define@term}
% \cs{hit@declare@class@key}\marg{名} 把一个键注册进 \texttt{hit/class} 族，处理器直接
% 调用同名命令。\cs{hit@define@term}\marg{名} 在此之上再造出那个同名命令：
% \cs{title-zh}\marg{值} 存进 \cs{hit@ctitle}，并且先置一次空值。
%    \begin{macrocode}
\cs_new_protected:Npn \hit@declare@class@key #1
  { \keys_define:nn { hit / class } { #1 .code:n = { \use:c {#1} {##1} } } }
\cs_new_protected:Npn \__hit_define_term:nn #1#2
  {
    \cs_gset_nopar:cpn {#1} ##1 { \cs_gset_nopar:cpn { hit@#2 } {##1} }
    \use:c {#1} { }
  }
\cs_new_protected:Npn \hit@define@term #1
  {
    \hit@declare@class@key {#1}
    \__hit_key_to_macro_name:n {#1}
    \exp_args:Nnx \__hit_define_term:nn {#1} { \l__hit_constant_name_str }
  }
%    \end{macrocode}
% \end{macro}
%
% \begin{macro}{\hit@define@legacy@terms}
% \changes{v3.2a}{2026/08/07}{元信息字段改成英文 kebab 名，旧名走废弃期}
% 元信息字段在 v3.1e 就是公开接口，改名后旧名要继续能用。
% \cs{hit@define@legacy@terms}\marg{旧名=新名,\ldots} 给每个旧名做三件事：把同名命令
% 指向新字段的设值器、把旧键注册进 \texttt{hit/class} 族、往改名表里记一条，
% 这样用户写旧名时会收到一次提示，告诉他新名叫什么。
%    \begin{macrocode}
\prop_new:N \l__hit_legacy_term_prop
\cs_new_protected:Npn \hit@define@legacy@term #1#2
  {
    \cs_gset_eq:cc {#1} {#2}
    \hit@declare@class@key {#1}
    \prop_gput:Nnn \g__hit_renamed_option_prop {#1} { info={#2=...} }
  }
\cs_new_protected:Npn \hit@define@legacy@terms #1
  {
    \prop_set_from_keyval:Nn \l__hit_legacy_term_prop {#1}
    \prop_map_function:NN \l__hit_legacy_term_prop \hit@define@legacy@term
  }
%    \end{macrocode}
% \end{macro}
%
% \begin{macro}{\hit@define@term@alias}
% \changes{v3.2a}{2026/08/11}{加 \cs{hit@define@term@alias}，两个名字都是正式写法}
% \cs{hit@define@term@alias}\marg{别名}\marg{正名} 也是把两个名字接到同一个设值器上，
% 但不进改名表，所以写哪个都不提示。跟 \cs{hit@define@legacy@term} 的区别在于
% 定位：那个是「旧名，该改了」，这个是「两个名字都对」。
% 学号写 \option{student-id} 还是 \option{student-number} 全看习惯，没有哪个更规范。
%    \begin{macrocode}
\cs_new_protected:Npn \hit@define@term@alias #1#2
  {
    \cs_gset_eq:cc {#1} {#2}
    \hit@declare@class@key {#1}
  }
%    \end{macrocode}
% \end{macro}
%
% \begin{macro}{\hit@define@date@term}
% \changes{v3.2a}{2026/08/13}{加 \cs{hit@define@date@term}，日期字段存原值、
%   取值时按样式与精度排}
% \cs{hit@define@date@term}\marg{键名}\marg{样式}\marg{精度} 与 \cs{hit@define@term}
% 并列，区别在取值那一头：原值存进 \cs{hit@}\meta{名}\cs{@raw}，
% \cs{hit@}\meta{名} 展开时才按样式与精度排。这样各处的排版代码一个字都不用改，
% 照旧写 \cs{hit@date@zh}。
%
% 精度可以是一段判断，用到时才求值——中文封面本科要年月日、硕博要年月，
% 是同一个字段两种精度，得等选项定下来才知道。
%
% 取值先经 \cs{exp\_args:NNv} 落进一个变量再传。直接把 \cs{use:c}\marg{名} 当参数
% 递过去，格式化那头拿到的是这几个记号本身，一看就不是 ISO，原样吐回来。
%    \begin{macrocode}
\cs_new_protected:Npn \__hit_define_date_term:nnnn #1#2#3#4
  {
    \cs_gset_nopar:cpn {#1} ##1 { \cs_gset_nopar:cpn { hit@#2@raw } {##1} }
    \use:c {#1} { }
    \cs_gset_protected_nopar:cpn { hit@#2 }
      {
        \exp_args:NNv \tl_set:Nn \l__hit_date_raw_tl { hit@#2@raw }
        \hit@date@use:nnnV {#1} {#3} {#4} \l__hit_date_raw_tl
      }
  }
\cs_generate_variant:Nn \__hit_define_date_term:nnnn { nxnn }
\cs_new_protected:Npn \hit@define@date@term #1#2#3
  {
    \hit@declare@class@key {#1}
    \__hit_key_to_macro_name:n {#1}
    \__hit_define_date_term:nxnn
      {#1} { \l__hit_constant_name_str } {#2} {#3}
  }
%    \end{macrocode}
% \end{macro}
%
% \begin{macro}{\hit@parse@keywords}
% \cs{hit@parse@keywords}\marg{字段名}\marg{分隔符宏名} 造出 \cs{keywords-zh}\marg{甲,乙,丙}，
% 把逗号表拼成 \cs{hit@ckeywords}，词与词之间插一个分隔符。分隔符是以 \cs{use:c} 的
% 形式存进去的，取值在排版时才解析，用户后改分隔符照样生效。分隔符的名字显式传进来，
% 不从字段名拼——两者的命名规律不一样。
%
% 同时另存一份 \cs{hit@}\meta{字段}\cs{@plain}：每项去掉首尾空格，词之间只有分隔符，
% 没有 \cs{ignorespaces}。这一份给 PDF 属性用——那条路上 \cs{ignorespaces} 不能执行，
% 屏蔽掉之后示例里 \texttt{\{\string\hologo{TeX}, \string\hologo{LaTeX}\}} 逗号后的空格就会漏进
% 元数据，排出来的版面却是没有空格的，两处对不上。拆分用 \cs{seq\_set\_split\_keep\_spaces:Nnn} 而不
% 是 \cs{clist\_map\_inline:nn}，因为后者会把每项前后的空格吃掉，而封面示例里写的是
% \texttt{\{\string\hologo{TeX}, \string\hologo{LaTeX}, CJK\}} 这种逗号后带空格的形式，吃掉就
% 与原先的 \cs{@for} 不一致了。
%    \begin{macrocode}
\seq_new:N \l__hit_keywords_seq
\tl_new:N \l__hit_keywords_item_tl
\cs_new_protected:Npn \__hit_parse_keywords:nnn #1#2#3
  {
    \cs_gset_nopar:cpn { hit@#3 } { }
    \cs_gset_nopar:cpn { hit@#3@plain } { }
    \cs_gset_protected_nopar:cpn {#1} ##1
      {
        \seq_set_split_keep_spaces:Nnn \l__hit_keywords_seq { , } {##1}
        \seq_map_inline:Nn \l__hit_keywords_seq
          {
            % 排版版的相邻项间不能夹宏(\ignorespaces 会把 xeCJK 的字符对
            % 拆成边界对,「！；」闭闭压缩失效——关键词行实测),改成存储时
            % 修剪条目两端空格,分隔符裸拼
            \tl_set:Nn \l__hit_keywords_item_tl {####1}
            \tl_trim_spaces:N \l__hit_keywords_item_tl
            \tl_if_empty:cF { hit@#3 }
              { \tl_gput_right:cx { hit@#3 } { \exp_not:v { hit@#2 } } }
            \tl_gput_right:cV { hit@#3 } \l__hit_keywords_item_tl
            \tl_if_empty:cF { hit@#3@plain }
              { \tl_gput_right:cn { hit@#3@plain } { \use:c { hit@#2 } } }
            \tl_gput_right:cx { hit@#3@plain } { \tl_trim_spaces:n {####1} }
          }
      }
  }
\cs_new_protected:Npn \hit@parse@keywords #1#2
  {
    \hit@declare@class@key {#1}
    \__hit_key_to_macro_name:n {#1}
    \exp_args:Nnnx \__hit_parse_keywords:nnn {#1} {#2} { \l__hit_constant_name_str }
  }
%    \end{macrocode}
% \end{macro}
%
% \begin{macro}{\hitsetup}
% 先占个位，各个类在自己的 meta 模块末尾用 \cs{RenewDocumentCommand} 接管。
%    \begin{macrocode}
\cs_new_protected:Npn \hitsetup #1 { }
%    \end{macrocode}
% \end{macro}
%
% \begin{macro}{\hit@keys@set:nn}
% \changes{v3.2a}{2026/08/13}{加 \cs{hit@keys@set:nn}，把键之间的空行剥掉再喂给
%   \cs{keys\_set:nn}；凡是接用户写的键值表的地方都改走它}
% \textbf{键之间夹了空行也得认。} 名单一条条写下来，中间空一行分组是很自然的写法，
% 但空行在键值表里是一个 \cs{par}，下一个键的名字就成了 \verb|\par secretary|，
% 对不上任何键。轻则丢一条记录只在日志里留行警告，重则被当成老写法去查。
%
% 所以喂给 \cs{keys\_set:nn} 之前先把顶层的 \cs{par} 剥掉。
% \cs{tl\_replace\_all:Nnn} 不进花括号，所以剥的只是键与键之间那些，
% 值里面的分段留得住——决议正文的空行是内容，不能动。
%
% 每一层都要剥。\cs{hitsetup} 那一层剥了，\texttt{defense~=~\{~\ldots~\}} 里面
% 是一个花括号组，没被碰到，所以子族转发时得再剥一次。凡是接用户写的键值表的
% 地方一律走这个函数，各层就都通了。
%    \begin{macrocode}
\tl_new:N \l__hit_keys_list_tl
\cs_new_protected:Npn \hit@keys@set:nn #1#2
  {
    \tl_set:Nn \l__hit_keys_list_tl {#2}
    \tl_replace_all:Nnn \l__hit_keys_list_tl { \par } { }
    \exp_args:NnV \keys_set:nn {#1} \l__hit_keys_list_tl
  }
%    \end{macrocode}
% \end{macro}
%
% 子族的兜底。键名先自己展开，靠 \cs{keys\_set:nn} 去展开 \cs{l\_keys\_key\_str} 不
% 可靠：载入的宏包集合一变就可能失效，换定理实现时踩过。
%    \begin{macrocode}
\cs_new_protected:Npn \__hit_legacy_key:n #1
  {
    \__hit_legacy_warn:
    \use:x
      { \exp_not:N \keys_set:nn { hit / class } { \l_keys_key_str = { \exp_not:n {#1} } } }
  }
\clist_map_inline:nn { info , style , toc , bib , const }
  { \keys_define:nn { hit / #1 } { unknown .code:n = { \__hit_legacy_key:n {##1} } } }
%    \end{macrocode}
%
% \changes{v3.2a}{2026/08/07}{\texttt{info} 与 \texttt{const} 两族加 \texttt{zh} 与
%   \texttt{en} 分组写法，把语言后缀提出来写一次}
% 中英文成对的键占了一小半，一个个写后缀有点啰嗦。\texttt{zh} 与 \texttt{en} 两个
% 分组把后缀提出来：
% \begin{latex}
% \hitsetup{ info = {
%   zh = { title = {中文题目}, author = {张三} },
%   en = { title = {English Title}, author = {Zhang San} },
%   student-id = {1180800000},
% } }
% \end{latex}
% 分组只是把里面每个键名接上后缀再转发，\texttt{zh~=~\{~title~=~\ldots~\}} 与
% \texttt{title-zh~=~\ldots} 完全等价。正式名仍然是带后缀的那个，报错与改名提示
% 里出现的都是它。没有配对的键（\texttt{student-id}、\texttt{secret-level} 这些）
% 照旧写在分组外面。
%
% \changes{v3.2a}{2026/08/09}{\texttt{zh} 与 \texttt{en} 改成真的子族，
%   斜杠路径写法跟着一起通}
% \texttt{zh} 与 \texttt{en} 是真的子族 \texttt{hit/info/zh}，接后缀这件事挂在子族的
% \texttt{unknown} 上，分组键自己只是一句 \cs{keys\_set:nn}。子族存在还带来一个用法：
% l3keys 按最后一个斜杠切键名，左边当族名右边当键名，所以斜杠路径也认。
% \begin{latex}
% \hitsetup{ info/zh/title = {中文题目}, info/title-en = {English Title} }
% \end{latex}
% 三种写法落到同一个键上，键名写错时报的也是同一条提示。
%
%    \begin{macrocode}
\cs_new_protected:Npn \hit@warn@renamed@environment #1#2
  { \__hit_warn_once:nnn {#1} { renamed-environment } {#2} }
\cs_new_protected:Npn \hit@warn@renamed@command #1#2
  { \__hit_warn_once:nnn {#1} { renamed-command } {#2} }
\cs_new_protected:Npn \__hit_language_key:nnnn #1#2#3#4
  {
    \keys_if_exist:nnTF { hit / #1 } { #3-#2 }
      { \keys_set:nn { hit / #1 } { #3-#2 = {#4} } }
      { \msg_warning:nnnn { hithesis } { key-not-in-language-group } {#3} {#2} }
  }
\cs_new_protected:Npn \__hit_language_unknown:nnn #1#2#3
  {
    \use:x
      {
        \exp_not:N \__hit_language_key:nnnn
          {#1} {#2} { \l_keys_key_str } { \exp_not:n {#3} }
      }
  }
\cs_new_protected:Npn \__hit_declare_language_group:nn #1#2
  {
    \keys_define:nn { hit / #1 / #2 }
      { unknown .code:n = { \__hit_language_unknown:nnn {#1} {#2} {##1} } }
    \keys_define:nn { hit / #1 }
      { #2 .code:n = { \hit@keys@set:nn { hit / #1 / #2 } {##1} } }
  }
\clist_map_inline:nn { info , const }
  {
    \clist_map_inline:Nn \c__hit_lang_clist
      { \__hit_declare_language_group:nn {#1} {##1} }
  }
%    \end{macrocode}
%
% 三态取值。选项可取 \texttt{auto}、\texttt{true}、\texttt{false}，也可以按学位分别
% 指定；\texttt{auto} 的含义由各选项自己定义，在用到的时候才解析。分学位映射里没
% 列到的学位保持 \texttt{auto}，不是 \texttt{false}。
%    \begin{macrocode}
\tl_new:N \l__hit_tristate_tl
\keys_define:nn { hit / degree-map }
  {
    bachelor .code:n = { \bool_if:NT \g__hit_bachelor_bool { \tl_set:Nn \l__hit_tristate_tl {#1} } } ,
    master   .code:n = { \bool_if:NT \g__hit_master_bool { \tl_set:Nn \l__hit_tristate_tl {#1} } } ,
    doctor   .code:n = { \bool_if:NT \g__hit_doctor_bool { \tl_set:Nn \l__hit_tristate_tl {#1} } } ,
    postdoc  .code:n = { \bool_if:NT \g__hit_postdoc_bool { \tl_set:Nn \l__hit_tristate_tl {#1} } } ,
  }
%    \end{macrocode}
% 设值。带等号的当作分学位映射，其余直接存。
%    \begin{macrocode}
\cs_new_protected:Npn \__hit_tristate_set:Nn #1#2
  {
    \tl_if_in:nnTF {#2} { = }
      {
        \tl_set:Nn \l__hit_tristate_tl { auto }
        \keys_set:nn { hit / degree-map } {#2}
        \tl_set_eq:NN #1 \l__hit_tristate_tl
      }
      { \tl_set:Nn #1 {#2} }
  }
%    \end{macrocode}
% 取值。\texttt{auto} 时执行第二个参数给出的布尔表达式来决定。
%    \begin{macrocode}
\cs_new_protected:Npn \__hit_tristate_if:NnTF #1#2#3#4
  {
    \str_case:VnF #1
      { { true } {#3} { false } {#4} }
      { \bool_if:nTF {#2} {#3} {#4} }
  }
\cs_new_protected:Npn \__hit_tristate_if:NnT #1#2#3
  { \__hit_tristate_if:NnTF #1 {#2} {#3} { } }
%    \end{macrocode}
%
% 造一个普通布尔键：\cs{\_\_hit\_boolean\_key:nnn}\marg{族}\marg{键名}\marg{标志位名}，
% 标志位是 \cs{g\_\_hit\_}\meta{名}\cs{\_bool}。新名设的是旧标志位本身，所以两套
% 名字天然一致。
%    \begin{macrocode}
\cs_new_protected:Npn \__hit_boolean_key:nnn #1#2#3
  {
    \use:x
      {
        \keys_define:nn { hit / #1 }
          {
            #2 .choice: ,
            #2 / true  .code:n = { \exp_not:N \bool_gset_true:c { g__hit_#3_bool } } ,
            #2 / false .code:n = { \exp_not:N \bool_gset_false:c { g__hit_#3_bool } } ,
            #2 / auto  .code:n = { } ,
            #2 .default:n = true ,
          }
      }
  }
\cs_new_protected:Npn \__hit_style_pair:nn #1#2 { \__hit_boolean_key:nnn { style } {#1} {#2} }
%    \end{macrocode}
%
% \begin{macro}{\hit@declare@info@field}
% 元信息字段注册进 \texttt{hit/info} 族。这批字段由 \cs{hit@define@term} 与
% \cs{hit@parse@keywords} 批量生成，每个都有一个同名命令（\cs{title-zh} 等），
% 键的处理器直接调用那个命令，所以三种写法结果一致：\cs{hitsetup} 的扁平写法、
% 分族写法、以及命令写法。名单由各个类传进来。
%    \begin{macrocode}
\clist_new:N \l__hit_info_clist
\prop_new:N \g__hit_info_prop
\cs_new_protected:Npn \hit@declare@info@field #1
  {
    \clist_set:Nn \l__hit_info_clist {#1}
    \clist_map_inline:Nn \l__hit_info_clist
      {
        \keys_define:nn { hit / info } { ##1 .code:n = { \use:c { ##1 } {####1} } }
        \prop_gput:Nnn \g__hit_info_prop {##1} { }
      }
  }
%    \end{macrocode}
% \end{macro}
%
% \begin{macro}{\hit@declare@constant,\hit@set@constant}
% 类文件里定义了一批可本地化的字符串：封面各字段的标签、声明与授权的正文、
% 签名栏文字等。各院系叫法不一，用户有改的需要，原先要改就得动生成物。
% 这里把它们做成键：
% \begin{latex}
% \hitsetup{ const = { supervisor-field-label-zh = {导师} } }
% \end{latex}
% 键名由宏名机械地得来：去掉 \texttt{hit@} 前缀，其余的 \texttt{@} 换成连字符，
% 所以 \cs{hit@ckeywords@separator} 对应 \texttt{keywords-separator-zh}。
%
% \cs{hit@set@constant} 把一张键值表灌进 \texttt{hit/const} 族。默认值与用户在导言区
% 写的 \cs{hitsetup}\{const=\{…\}\} 走的是同一条路，区别只是先后：类文件末尾先设
% 默认值，导言区后设的覆盖它。表是在 expl3 之外记号化的，取值里的空格照原样保留。
%
% \changes{v3.2a}{2026/08/07}{常量取值里可以用 \texttt{<键名!>} 引用别的常量，
%   不必再写内部宏名}
% \changes{v3.2a}{2026/08/07}{删掉 \texttt{cxueweishort}：它只为拼
%   \texttt{degree-level-zh} 与 \texttt{author-field-label-zh} 存在，两个都能从 \texttt{degree-level-zh}
%   直接写出来。本科分支从没设过它，那一处引用的是未定义的宏}
% 有些取值要拼进别的常量。取值里写 \texttt{<键名>} 就行，尖括号里是手册上那个键名：
% \begin{latex}
% \hitsetup{ const = { author-field-label-zh = {<cxueweishort>士研究生} } }
% \end{latex}
% 替换发生在设值时（\texttt{<键名>} 换成对应的控制序列），取值仍发生在排版时，所以
% 被引用的常量之后再改也跟得上。尖括号里只认小写字母与连字符，取值里真正的
% \texttt{<} 不会被误伤；直接写 \cs{hit@}\meta{名} 的老写法照旧可用。
%
% 这样写还避开一个坑：\cs{hit@cxueweishort} 是控制字，后面必须留一个空格来终止它，
% 那个空格不会排出来。照抄的人容易以为空格是版式的一部分。\texttt{<键名>} 后面没有
% 控制字，空格就是空格。
%    \begin{macrocode}
\str_new:N \l__hit_constant_name_str
\tl_new:N \l__hit_constant_value_tl
\tl_new:N \l__hit_reference_tl
\seq_new:N \l__hit_reference_seq
\cs_generate_variant:Nn \tl_replace_all:Nnn { Nnx }
\cs_generate_variant:Nn \regex_extract_all:nnN { nVN }
\cs_new_protected:Npn \__hit_key_to_macro_name:n #1
  {
    \str_set:Nn \l__hit_constant_name_str {#1}
    \str_replace_all:Nnn \l__hit_constant_name_str { - } { @ }
  }
\cs_generate_variant:Nn \__hit_key_to_macro_name:n { V }
\cs_new_protected:Npn \__hit_expand_references:N #1
  {
    \regex_extract_all:nVN { < [a-z\-]+ > } #1 \l__hit_reference_seq
    \seq_remove_duplicates:N \l__hit_reference_seq
    \seq_map_inline:Nn \l__hit_reference_seq
      {
        \tl_set:Nn \l__hit_reference_tl {##1}
        \tl_remove_all:Nn \l__hit_reference_tl { < }
        \tl_remove_all:Nn \l__hit_reference_tl { > }
        \__hit_key_to_macro_name:V \l__hit_reference_tl
        \tl_replace_all:Nnx #1 {##1}
          { \exp_not:c { hit@ \l__hit_constant_name_str } }
      }
  }
\cs_new_protected:Npn \__hit_constant_set:n #1
  {
    \tl_set:Nn \l__hit_constant_value_tl {#1}
    \__hit_expand_references:N \l__hit_constant_value_tl
    \__hit_key_to_macro_name:V \l_keys_key_str
    \cs_gset:cpx { hit@ \l__hit_constant_name_str }
      { \exp_not:V \l__hit_constant_value_tl }
  }
\cs_new_protected:Npn \hit@declare@constant #1
  {
    \clist_map_inline:nn {#1}
      { \keys_define:nn { hit / const } { ##1 .code:n = { \__hit_constant_set:n {####1} } } }
  }
%    \end{macrocode}
%
% \changes{v3.2a}{2026/08/08}{常量可以直接驱动既有命令名或转发给 \pkg{ctex}，
%   配置里那批 \cs{cs\_new:Npn} \cs{equationname} 与 \cs{ctexset} 因此不必再写在
%   配置里——配置只出取值，造轮子的部分都在这一层}
% 有两类可本地化的词不走 \cs{hit@}\meta{名} 这个存法：一类是 \cs{equationname}、
% \cs{cabstractcname} 这种早就对外可见的名字，改名会破坏接口；一类是
% \cs{contentsname}、\cs{figurename} 这种归 \pkg{ctex} 管的名字，得走 \cs{ctexset}。
%
% \cs{hit@declare@constant@as}\marg{键名列表} 声明的键，值直接写进同名命令
% （\texttt{equationname} 写 \cs{equationname}，不是 \cs{hit@equationname}）。
% \cs{hit@declare@constant@ctex}\marg{键名=ctex键,\ldots} 声明的键，值转给
% \cs{ctexset}；两个名字常常一样，不一样的写等号（\texttt{chapter-name=chapter/name}）。
%
% 两者都照旧认取值里的 \texttt{<键名>} 引用。
%    \begin{macrocode}
%    \end{macrocode}
%
% \changes{v3.2a}{2026/08/11}{加 \cs{hit@after@number:nnn}，让「编号之后排什么」
%   这批常量取空值时退回各处原有的取值}
% \cs{hit@after@number:nnn}\marg{层级}\marg{语言}\marg{兜底}
% 取 \option{after-}\meta{层级}\option{-number-}\meta{语言} 这个常量，取值为空就用兜底。
%
% 「编号之后排什么」这批常量的默认值不能写成一个数：章号之后排多宽，本部硕士与
% 深圳本科是一个全角，其余是半个；报告类还要看开题还是中期。这些判断本来就写在
% 各个 \cs{ctexset} 块里，一个块一份。所以常量默认留空，各块把自己原有的取值当
% 兜底传进来，用户没设就跟以前一模一样，设了就一处覆盖全部。
%
% 空不空不能用 \cs{tl\_if\_empty:cTF} 判。常量是 \cs{cs\_gset:cpx} 建的，那是个
% \cs{long} 宏；\cs{tl\_if\_empty:N} 拿 \cs{c\_empty\_tl} 比意义，而
% \cs{c\_empty\_tl} 不是 \cs{long}，空常量会被判成非空，接着 \cs{use:c} 取到那个
% 空宏，间距静悄悄地消失。踩过一次，症状是 45 个变体全变、页数一页不变。
% 所以另造一个同样由 \cs{cs\_new:Npn} 建的空宏来比。
%
% 常量不存在时也走兜底：默认值虽然在 \file{hit-defaults.dtx} 里都设了，
% 少一条不该让间距无声消失。
%    \begin{macrocode}
\cs_new:Npn \hit@after@number@empty: { }
\cs_new:Npn \hit@after@number:nnn #1#2#3
  {
    \cs_if_exist:cTF { hit@after@#1@number@#2 }
      {
        \cs_if_eq:cNTF { hit@after@#1@number@#2 } \hit@after@number@empty:
          {#3}
          { \use:c { hit@after@#1@number@#2 } }
      }
      {#3}
  }
\tl_new:N \l__hit_constant_target_tl
\prop_new:N \g__hit_constant_target_prop
\cs_new_protected:Npn \__hit_constant_target:
  {
    \prop_get:NVN \g__hit_constant_target_prop \l_keys_key_str
      \l__hit_constant_target_tl
  }
\cs_new_protected:Npn \__hit_constant_set_as:n #1
  {
    \tl_set:Nn \l__hit_constant_value_tl {#1}
    \__hit_expand_references:N \l__hit_constant_value_tl
    \__hit_constant_target:
    \cs_gset:cpx { \l__hit_constant_target_tl }
      { \exp_not:V \l__hit_constant_value_tl }
  }
\cs_new_protected:Npn \__hit_constant_set_ctex:n #1
  {
    \tl_set:Nn \l__hit_constant_value_tl {#1}
    \__hit_expand_references:N \l__hit_constant_value_tl
    \__hit_constant_target:
    \exp_args:Nx \ctexset
      { \l__hit_constant_target_tl = { \exp_not:V \l__hit_constant_value_tl } }
  }
\cs_new_protected:Npn \hit@declare@constant@as #1
  {
    \clist_map_inline:nn {#1}
      {
        \prop_gput:Nnn \g__hit_constant_target_prop {##1} {##1}
        \keys_define:nn { hit / const }
          { ##1 .code:n = { \__hit_constant_set_as:n {####1} } }
      }
  }
\cs_new_protected:Npn \hit@declare@constant@ctex #1
  { \clist_map_inline:nn {#1} { \__hit_declare_constant_ctex:w ##1 = ##1 = \q_stop } }
\tl_new:N \l__hit_constant_key_tl
\cs_generate_variant:Nn \prop_gput:Nnn { NVV }
\cs_new_protected:Npn \__hit_declare_constant_ctex:w #1 = #2 = #3 \q_stop
  {
    \tl_set:Nx \l__hit_constant_key_tl { \tl_trim_spaces:n {#1} }
    \tl_set:Nx \l__hit_constant_target_tl { \tl_trim_spaces:n {#2} }
    \prop_gput:NVV \g__hit_constant_target_prop
      \l__hit_constant_key_tl \l__hit_constant_target_tl
    \exp_args:Nnx \keys_define:nn { hit / const }
      {
        \l__hit_constant_key_tl .code:n =
          { \exp_not:N \__hit_constant_set_ctex:n { \exp_not:n {##1} } }
      }
  }
%    \end{macrocode}
%
%    \begin{macrocode}
\cs_new_protected:Npn \hit@set@constant #1 { \hit@keys@set:nn { hit / const } {#1} }
%    \end{macrocode}
% \end{macro}
%
% 兼容期的提示文案。\texttt{v4} 是兼容期结束后的第一个架构大版本。
%    \begin{macrocode}
%    \end{macrocode}
%
% 报过的名字记下来，同一个名字只报一次。元信息字段合并成一条，用第一个碰到的
% 字段名举例。旧名到新名的映射表 \cs{g\_\_hit\_renamed\_option\_prop} 由各个类填。
%
% \changes{v3.2a}{2026/08/10}{改名映射分成两张表。\cs{g\_\_hit\_renamed\_option\_prop} 装
%   排版选项，写进 \cs{documentclass} 会另报一条；\cs{g\_\_hit\_renamed\_load\_option\_prop}
%   装加载期选项，写进 \cs{documentclass} 本来就对，只报改名}
% 两张改名表。排版选项写在 \cs{documentclass} 里要提示改用 \cs{hitsetup}；
% 加载期选项（字体、校区、学位这类）必须写在 \cs{documentclass} 里，只提示改名。
%
% \changes{v3.2a}{2026/08/11}{加废弃表 \cs{g\_\_hit\_removed\_option\_prop}，
%   装没有新名可改的选项}
% 第三张是废弃表，装的是不再有对应写法的选项。它排在两张改名表之前查：
% 一个键要么改了名要么废了，不会两者都是，先查哪张只影响没写对时的提示，
% 而废弃的那条话说得更具体。表里存的是「为什么废了」，直接进提示。
%    \begin{macrocode}
\prop_new:N \g__hit_removed_option_prop
\prop_new:N \g__hit_renamed_option_prop
\prop_new:N \g__hit_split_option_prop
\prop_new:N \g__hit_renamed_load_option_prop
\prop_new:N \g__hit_warned_options_prop
\tl_new:N \l__hit_warn_message_tl
\cs_generate_variant:Nn \prop_get:NnNTF { NV }
\cs_generate_variant:Nn \msg_warning:nnn { nnV }
\cs_new_protected:Npn \__hit_warn_once:nnn #1#2#3
  {
    \prop_if_in:NnF \g__hit_warned_options_prop {#1}
      {
        \prop_gput:Nnn \g__hit_warned_options_prop {#1} { }
        \msg_warning:nnnn { hithesis } {#2} {#1} {#3}
      }
  }
\cs_generate_variant:Nn \__hit_warn_once:nnn { VnV }
\bool_new:N \g__hit_flat_info_warned_bool
\cs_new_protected:Npn \__hit_legacy_warn:
  {
    \prop_get:NVNTF \g__hit_removed_option_prop \l_keys_key_str \l__hit_warn_message_tl
      {
        \__hit_warn_once:VnV \l_keys_key_str
          { removed-option } \l__hit_warn_message_tl
      }
      {
    \prop_get:NVNTF \g__hit_split_option_prop \l_keys_key_str \l__hit_warn_message_tl
      {
        \__hit_warn_once:VnV \l_keys_key_str
          { split-option } \l__hit_warn_message_tl
      }
      {
    \prop_get:NVNTF \g__hit_renamed_option_prop \l_keys_key_str \l__hit_warn_message_tl
      {
        \__hit_warn_once:VnV \l_keys_key_str
          { renamed-option } \l__hit_warn_message_tl
      }
      {
    \prop_get:NVNTF \g__hit_renamed_load_option_prop \l_keys_key_str \l__hit_warn_message_tl
      {
        \__hit_warn_once:VnV \l_keys_key_str
          { renamed-load-option } \l__hit_warn_message_tl
      }
      {
        \bool_lazy_and:nnT
          { ! \bool_if_p:N \g__hit_flat_info_warned_bool }
          { \prop_if_in_p:NV \g__hit_info_prop \l_keys_key_str }
          {
            \bool_gset_true:N \g__hit_flat_info_warned_bool
            \msg_warning:nnV { hithesis } { flat-info } \l_keys_key_str
          }
      }
      }
      }
      }
  }
%    \end{macrocode}
%
% \cs{documentclass} 的可选参数里出现排版选项时提示一次。\cs{@classoptionslist}
% 在选项处理完之后仍然保留，所以放到这里扫描。每扫到一个选项名，名字存在
% \cs{l\_\_hit\_warn\_new\_name\_tl} 里，接着执行一次 \cs{\_\_hit\_class\_option\_extra:}，类自己
% 要额外查什么就重定义它，默认什么都不做。
%    \begin{macrocode}
\clist_new:N \l__hit_class_options_clist
\tl_new:N \l__hit_warn_new_name_tl
\cs_generate_variant:Nn \prop_get:NnNT { NV }
\cs_generate_variant:Nn \__hit_warn_once:nnn { VnV }
\cs_new_protected:Npn \__hit_class_option_extra: { }
\cs_new:Npn \__hit_option_name:w #1 = #2 \q_stop { \tl_trim_spaces:n {#1} }
\cs_new_protected:Npn \__hit_check_class_options:
  {
    \cs_if_exist:cT { @classoptionslist }
      {
        \exp_args:NNv \clist_set:Nn \l__hit_class_options_clist { @classoptionslist }
        \clist_map_inline:Nn \l__hit_class_options_clist
          {
            \tl_set:Nx \l__hit_warn_message_tl { \__hit_option_name:w ##1 = \q_stop }
            \prop_if_in:NVT \g__hit_renamed_option_prop \l__hit_warn_message_tl
              {
                \exp_args:NV \__hit_warn_once:nnn \l__hit_warn_message_tl
                  { style-in-classoptions } { }
              }
            \prop_get:NVNT \g__hit_split_option_prop \l__hit_warn_message_tl \l__hit_warn_new_name_tl
              {
                \exp_args:NVnV \__hit_warn_once:nnn \l__hit_warn_message_tl
                  { split-option } \l__hit_warn_new_name_tl
              }
            \prop_get:NVNT \g__hit_renamed_load_option_prop \l__hit_warn_message_tl \l__hit_warn_new_name_tl
              {
                \exp_args:NVnV \__hit_warn_once:nnn \l__hit_warn_message_tl
                  { renamed-load-option } \l__hit_warn_new_name_tl
              }
            \__hit_class_option_extra:
          }
      }
  }
%    \end{macrocode}
%
%    \begin{macrocode}
%</shared-key-engine>
%    \end{macrocode}
%
% \Finale
